\documentclass[../../main.tex]{subfiles}
\begin{document}

\begin{center}
  {\fontsize{48pt}{58pt}\selectfont \textbf{1971}\normalsize\textbf{年度}}
\end{center}

\begin{table}[htb]
  \centering
  \caption{出題テーマと難易度}
  \begin{tabular}{|c|p{9cm}|c|}
    \hline
    \textbf{問題番号} & \textbf{テーマ} & \textbf{難易度} \\
    \hline
    第1問 & 正方形に含まれる格子点 & \\
    \hline
    第2問 & 共役複素数を含む一次方程式の解の存在 & \\
    \hline
    第3問 & 二次関数の区間最小値の最大化 & \\
    \hline
    第4問 & 絶対値を含む曲線で囲まれた回転体 & \\
    \hline
    第5問 & 正多角形と連続する扇形の面積 & \\
    \hline
    第6問 & 円周上の点と直線上の点の共線時刻 & \\
    \hline
  \end{tabular}
\end{table}

\begin{table}[htb]
  \centering
  \caption{解答}
  \renewcommand{\arraystretch}{1.4}
  \begin{tabular}{|c|c|p{9cm}|}
    \hline
    \textbf{問題番号} & \textbf{小問} & \textbf{答え} \\
    \hline
    第1問 & (1) & 軸に平行な一辺$1$の正方形は格子点を含む \\
    \cline{2-3}
    & (2) & 一辺$\sqrt{2}$の正方形は必ず格子点を含む \\
    \hline
    第2問 & -- & 解が存在する必要十分条件は$\alpha\bar{\beta}=\beta$（証明問題） \\
    \hline
    第3問 & -- & $(a,b)=(-1,\dfrac{3}{2})$のとき最大値$\dfrac{5}{4}$ \\
    \hline
    第4問 & -- & $V=\dfrac{4\pi}{3}(5\sqrt{5}-2)$ \\
    \hline
    第5問 & (1) & $S_n=\dfrac{nr^2}{2}\sin\dfrac{2\pi}{n}+\dfrac{2\pi r^2}{3}(n+1)(2n+1)\sin^2\dfrac{\pi}{n}$ \\
    \cline{2-3}
    & (2) & $\displaystyle\lim_{n\to\infty}S_n=\pi r^2+\dfrac{4}{3}\pi^3r^2$ \\
    \hline
    第6問 & (1) & 各区間$(n,n+1)$に共線時刻はただ1つ \\
    \cline{2-3}
    & (2) & $\displaystyle\lim_{n\to\infty}(t_n-n)=\dfrac{1}{2}$ \\
    \hline
  \end{tabular}
\end{table}

\paragraph{試験形態}
6問．試験時間は要確認．

\end{document}
